% !TeX root = ../../main.tex
\section{\bilingual{Limit Theorem for Log-Scaled Space-Time \texorpdfstring{$N$}{N}-BRW}{对数缩放后一类 \texorpdfstring{\!$N$}{N}-分支随机游走的极限定理}}

    \bilingual{
        A\"id\'ekon and Hu \mycite{aïdékon2025averageweightpercolationcompletegraph} proved a limit theorem for a special branching law $\Xi$ under which the $N$-branching random walk is defined. This branching law is given by
        \begin{equation*}
            \Xi=\sum_{i=1}^N\delta_{eT_i-1},
        \end{equation*}
        where $0<T_1<T_2<\cdots$ is a Poisson point process on $\mathbb R_+$ with parameter $1$. Their limit theorem states that, for the $N$-branching random walk starting with $\lfloor N^a\rfloor$ particles at position $0$ at time $0$ ($a\in[0,1]$), under the space-time scaling of log-cubed in time and single logarithm in space, the position of the leftmost particle at each time slice $t>0$ converges in $\mathrm L^2$ to a position on a straight line $-a+\frac{\pi^2}2t$ in this log-scaled space-time as $N\to\infty$, where $\frac{\pi^2}2$ comes from the application of \nameref{theorem:mogulskii}. In this thesis, following their method, we will generalize their result to more general assumptions in \Cref{proposition:3.1} in Section \ref{section:main-results}. Moreover, in the final conclusion about the order of the asymptotic speed of the leftmost particle of $N$-branching random walk (\Cref{theorem:speed}), this generalized result will be applied.
    }{
        A\"id\'ekon 与 Hu \mycite{aïdékon2025averageweightpercolationcompletegraph} 为一种特殊分支律 $\Xi$ 下的 $N$-分支随机游走证得了一个极限定理. 该分支律为
        \begin{equation*}
            \Xi=\sum_{i=1}^N\delta_{eT_i-1},
        \end{equation*}
        其中 $0<T_1<T_2<\cdots$ 是一个 $\mathbb R_+$ 上以 $1$ 为参数的 Poisson 点过程. 他们的极限定理表明,~ 对于零时刻从位置 $0$ 以 $\lfloor N^a\rfloor$ 个粒子出发的 $N$-分支随机游走 ($a\in[0,1]$),~ 在三次对数的时间与一次对数的空间下,~ 每一个时间切片 $t>0$ 处的最左侧粒子的位置在 $N\to\infty$ 下的 $\mathrm L^2$-极限位置落于该对数缩放时空下的一条直线$-a+\frac{\pi^2}2t$上该时间切片处,~ 其中 $\frac{\pi^2}2$ 来自于对\nameref{theorem:mogulskii}的应用. ~本文将在第 \ref{section:main-results} 节中的\Cref{proposition:3.1},~ 遵照他们的方法,~ 推广他们的结果至更一般的假设条件. ~并且,~ 在最后的关于 $N$-分支随机游走最左侧粒子的渐近速率阶数的结论 (\Cref{theorem:speed}) 中,~ 这一推广的结果将被应用.
    }
